\documentclass[12pt]{article}
\usepackage[left=2cm,right=2cm,top=2cm,bottom=2cm]{geometry}
\usepackage{wrapfig}
\usepackage{booktabs}
\usepackage{graphicx}
\usepackage{amsmath}
\usepackage{siunitx} % Required for alignment
\usepackage{subfigure}
\usepackage{multirow}
\usepackage{rotating}
\usepackage[T2A]{fontenc}
\usepackage[english, russian]{babel}
\usepackage{caption}
\usepackage{hyperref}
\usepackage{float}

\hypersetup{
    colorlinks=true,
    linkcolor=red,
    urlcolor=magenta,
    }
\graphicspath{{pictures/}}


\title{\begin{center}Лабораторная работа №1.3.3\end{center}
Определение вязкости воздуха по скорости течения через тонкие трубки}
\author{Балдин Виктор Б01-303}
\date{\today}

\begin{document}
\maketitle

\paragraph*{\textbf{Цель работы:}} экспериментально исследовать свойства течения газов по тонким трубкам при
различных числах Рейнольдса; выявить область применимости закона Пуазейля и с его помощью
определить коэффициент вязкости воздуха.
\paragraph*{\textbf{В работе используются:}} система подачи воздуха (компрессор, поводящие трубки); газовый
счетчик барабанного типа; спиртовой микроманометр с регулируемым наклоном; набор трубок
различного диаметра с выходами для подсоединения микроманометра; секундомер.

\section{Теоретическая часть}
	Характер движения газа по трубке определяется числом Рейнольдса:
	\begin{equation}
		Re = \frac{u r \rho}{\eta},
		\label{eq:Re}
	\end{equation}
	где $u$ - скорость потока, $r$ - радиус трубки, $\rho$ - плотность жидкости, $\eta$ - вязкость. Переход от ламинарного движения к турбулентному: $Re \sim 1000$.

	Для ламинарного течения при постоянном удельном объеме верна формула Пуазейля:
	\begin{equation}
		Q_V = \frac{\pi r^4}{8 l \eta}(P_1 - P_2),
		\label{eq:Pu}
	\end{equation}
	где $P_1 - P_2$ - разность давлений в двух сечениях, расстояние между которыми - $l$. Формула позволяет определить вязкость по расходу $Q_V$.

	Ламинарное течение газа устанавливается на расстоянии
	\begin{equation}
		a \sim 0.2 r \cdot Re.
		\label{eq:Aa}
	\end{equation}
	Градиент давления на участке с турбулентным течением больше, чем на участке с ламинарным, что позволяет разделить их экспериментально.

\section{Ход работы}
\begin{enumerate}
    \item Используем 2 трубки, диаметры и оптимальные расстояния, вычисленные по формуле \ref{eq:Aa},
    представлены в таблице \ref{tab:trub}.
    \begin{table}[h!]
        \centering
        \begin{tabular}{|c|c|c|}
        \toprule
        № трубки & $d$, мм & $a$, мм\\
        \midrule
        1    & $5.25 \pm 0.05$ & $525 \pm 5$\\
        2    & $3.90 \pm 0.05$ & $390 \pm 5$\\
        \bottomrule
        \end{tabular}
        \caption{Используемые трубки}
        \label{tab:trub}
    \end{table}

    \item Установим вывод манометра на оптимальном расстоянии от начала трубок и проведем
    измерения зависимости $\Delta P(Q)$. Множитель манометра установим $k = 0.2$. Для трубки
    №1 результаты представлены в таблице \ref{tab:wide}, для трубки №2 -- в таблице \ref{tab:thin}.
    \begin{table}[h!]
        \centering
        \begin{tabular}{|r|r|r|}
        \toprule
          $P$, мм &    $t$, c &     $Q$, см$^3$/c \\
        \midrule
         21 & 25.0 &  40.0 \\
         38 & 16.0 &  62.5 \\
         54 & 12.0 &  83.3 \\
         86 &  8.2 & 122.0 \\
        100 &  7.5 & 133.3 \\
         30 & 19.4 &  51.5 \\
         45 & 14.0 &  71.4 \\
         61 & 10.9 &  91.7 \\
         70 &  9.7 & 103.1 \\
         74 &  9.2 & 108.7 \\
         77 &  8.9 & 112.4 \\
        \bottomrule
        \end{tabular}
        \caption{$\Delta P(Q)$ для трубки №1}
        \label{tab:wide}
        \end{table}

        \begin{table}[h!]
        \centering
        \begin{tabular}{|r|r|r|}
        \toprule
          $P$, мм &    $t$, c &     $Q$, см$^3$/c \\
        \midrule
         19 & 48.5 &  20.6 \\
         30 & 33.0 &  30.3 \\
         40 & 25.0 &  40.0 \\
         50 & 20.6 &  48.5 \\
         60 & 18.1 &  55.2 \\
         70 & 16.0 &  62.5 \\
         80 & 14.0 &  71.4 \\
         91 & 12.6 &  79.4 \\
        100 & 11.7 &  85.5 \\
        110 & 11.0 &  90.9 \\
        106 & 11.2 &  89.3 \\
        126 & 10.1 &  99.0 \\
        140 &  9.6 & 104.2 \\
        \bottomrule
        \end{tabular}
        \caption{$\Delta P(Q)$ для трубки №2}
        \label{tab:thin}
        \end{table}

    \item Выделим линейные участки зависимости и построим на них наилучшие прямые,
    проходящие через 0, по МНК.
\end{enumerate}

\end{document}
